\section{ContextCoherenceBench}
\label{sec:benchmark}

To enable replication and follow-on work we release
\textbf{ContextCoherenceBench}, a public benchmark targeting the three
empirical phenomena reported in this paper. The release is bundled by
\texttt{scripts/package\_benchmark.py} and shipped to
HuggingFace Datasets and a companion GitHub repository (URLs withheld for
double-blind review; both are tagged \texttt{bench-v1} at the same git
SHA).

\subsection{Tasks}

\begin{description}
  \item[\texttt{adversarial\_coherence}] 40 hand-authored cases
        (10 \textsc{disjoint} / 10 \textsc{contradict} / 10
        \textsc{drift} / 10 \textsc{control}). Target metric: per-signal
        detection AUC against the matched control set, scored under the
        preregistered plan in §\ref{sec:adv_plan}. Designed to
        dissociate which coherence signal catches which fragmentation
        pattern.
  \item[\texttt{coherence\_paradox}] the multi-dataset query set
        (§\ref{sec:multidataset}) with per-query measurements for five
        conditions (baseline, hcpc\_v1, hcpc\_v2, crag, selfrag). Target
        metrics: \texttt{paradox\_drop} and \texttt{v2\_recovery}, under
        the reporting plan in §\ref{sec:multidataset_plan}. Lets
        researchers test whether their own retrievers / generators
        reproduce the paradox.
  \item[\texttt{human\_faithfulness}] 500 double-annotated items with
        binary faithfulness, 1--5 correctness, 1--5 helpfulness, and
        Fleiss' $\kappa$. Target metric: NLI $\leftrightarrow$ human
        Spearman correlation, under the reporting plan in
        §\ref{sec:human_plan}. Lets follow-on work calibrate new
        faithfulness scorers against the same human ratings.
\end{description}

\subsection{Layout}

The release directory mirrors the HuggingFace dataset spec:

\begin{verbatim}
release/context_coherence_bench_v1/
  README.md                  (HF dataset card)
  CITATION.bib  LICENSE
  metadata.json              (provenance + sha256 per file)
  adversarial/
    {disjoint,contradict,drift,control}.jsonl
  coherence_paradox/
    per_query.csv
    summary.csv
  human_faithfulness/
    aggregated.csv
    manifest.csv
\end{verbatim}

The hand-authored adversarial cases are released under \textbf{CC-BY-4.0}.
Per-dataset content inherits the upstream license recorded in
\texttt{metadata.json}; the human-annotation ratings are
\textbf{CC-BY-4.0} with annotator IDs redacted to opaque \texttt{A<n>}
strings.

\subsection{Intended use}

We see three productive uses for the benchmark:

\begin{enumerate}
  \item \textbf{Coherence-signal research.} The
        \texttt{adversarial\_coherence} subset isolates failure modes
        cleanly enough that a new signal proposal can be evaluated on
        category-stratified AUC rather than only aggregate detection.
  \item \textbf{Selective-refinement comparison.} The
        \texttt{coherence\_paradox} subset already reports five
        conditions on the same queries; a new retrieval method (e.g., a
        sixth condition) can be added without re-running ours.
  \item \textbf{Faithfulness-scorer calibration.} The
        \texttt{human\_faithfulness} subset gives any new
        automatic-faithfulness model a fixed reference set to validate
        against, including the paradox-survival check.
\end{enumerate}

We \emph{do not} intend the benchmark as a leaderboard. The metric the
paper centers on (\texttt{paradox\_drop}) is a diagnostic, not a
performance score; ranking systems by it would invert the intended use.
The release ships scripts but no scoreboard.

\subsection{Reproduction commands}

The companion repository ships entry points matching each subset:

\begin{verbatim}
python3 experiments/run_adversarial_coherence.py
python3 experiments/run_multidataset_validation.py
python3 experiments/run_headtohead_comparison.py
python3 experiments/run_mechanistic_analysis.py
python3 experiments/prepare_prolific_study.py
python3 experiments/analyze_prolific_results.py
python3 scripts/package_benchmark.py
\end{verbatim}

Each script writes a self-contained results directory under
\texttt{results/} that the packager copies into the release tree. Re-running
any single experiment refreshes only its own subset, so partial reruns
remain coherent.

\subsection{Permanent identifier}

The benchmark bundle is archived on Zenodo with a permanent DOI:
\href{https://doi.org/10.5281/zenodo.19757291}
     {\texttt{10.5281/zenodo.19757291}}.
The DOI resolves to the v$2.0.0$ release tarball
(\texttt{context\_coherence\_bench\_v1.tar.gz}) and includes the
\texttt{CITATION.bib} entry below for citing this artifact independently
of the paper:

\begin{verbatim}
@dataset{maganti2026ccbench,
  title     = {ContextCoherenceBench: a benchmark for evaluating
               context coherence in retrieval-augmented generation},
  author    = {Maganti, Saket},
  year      = {2026},
  publisher = {Zenodo},
  version   = {v2.0.0},
  doi       = {10.5281/zenodo.19757291},
}
\end{verbatim}
